\section{Introduction}
Reactive systems must continuously respond to an environment, often under
strict timing and safety constraints. Programming such systems is difficult
because concurrency, external events, and state updates interact in ways that
can break determinism. Ad-hoc callback-based designs are common, but they make
control flow implicit and complicate reasoning about ordering, causality, and
reproducibility.

Synchronous reactive programming addresses these challenges by discretizing
time into logical instants and by defining what is observable at instant
boundaries. A classic distinction is between control-flow and dataflow
approaches: Esterel represents the control-flow lineage, while Lustre
represents the dataflow lineage \cite{esterel1992,lustre1991}. In both cases,
logical instants enable deterministic behavior and rigorous analyses of
causality and scheduling. Survey discussions situate these languages within a
shared semantic framework \cite{benveniste2003}.

These synchronous languages achieve high assurance by imposing restrictions
that permit static scheduling and causality checks. The restrictions make it
possible to decide properties such as determinism and absence/presence
consistency, but they also limit expressiveness when integrated with richer
general-purpose languages.

The key semantic question is when signal absence can be observed. In Esterel,
absence is decided within the instant, enabling static scheduling but
restricting dynamic behavior \cite{esterel1992}. Boussinot's FairThreads model
defers absence observation to instant end, allowing more flexible composition
while preserving determinism \cite{fairthreads2003}. From a programmer's
perspective, each instant follows a two-phase discipline: react while tasks run
and signals may become present, then finalize when the instant quiesces and
absence becomes observable.
This distinction is central because it determines which control patterns are
expressible without sacrificing determinism.
Subsequent work revisits the synchronous model inside richer language settings.
On the control-flow side, ReactiveML adapts Esterel-style signals and the
synchronous hypothesis to a higher-order language \cite{reactiveml2005}. On the
dataflow side, FRP explores declarative, time-varying values with different
operational priorities \cite{frp1997}. 

This paper introduces \tempo{}: a library that revisits ReactiveML's control-flow
model using OCaml 5.3 algebraic effects as its implementation mechanism. \tempo{}
follows the delayed-absence discipline introduced by Boussinot's FairThreads,
which ReactiveML also adopts, and brings this approach to a modern OCaml
runtime without modifying the compiler \cite{fairthreads2003,reactiveml2005}.
\tempo{} re-implements the core control constructs (signals, guards, weak
preemption) as a pure library. It executes programs in discrete logical
instants: signals can be emitted and tasks can run within an instant, but the
\emph{absence} of a signal is only observable at the end of the instant. This
article presents \tempo{} as a reactive control runtime and explains its design
choices and effect-based implementation in OCaml 5.3. It does not attempt a
formal semantics or a quantitative evaluation, and focuses instead on the
design rationale, API, and implementation details needed to motivate future
formalization. Determinism is discussed under standard purity assumptions:
programs should avoid shared mutable state whose outcomes depend on
intra-instant scheduling, and aggregate combine functions should be
order-insensitive. These assumptions match the synchronous control model and
are stated explicitly in the core constructs section.
\begin{itemize}
\item A library-level formulation of synchronous reactive control in OCaml,
  positioned within the Esterel/FairThreads/ReactiveML lineage.
\item An implementation account of how algebraic effects realize signals,
  guards, suspension, and weak preemption.
\item End-to-end examples that illustrate standard synchronous control idioms
  using the \tempo{} API.
\end{itemize}
